\section{Experiments and Results}
\label{sec:experiments}

Following the comparative methodology established in compact vision architecture studies \cite{16}, we evaluate both overparameterized baselines and compact models constrained to a strict sub-1M parameter budget ($<10^6$ parameters). Overparameterized baselines comprise widely adopted backbones across computer vision and medical neuro-imaging: ResNet-50 and ResNet-18 \cite{10}, VGG-16 \cite{29}, EfficientNet-B0 \cite{30}, and MobileNet-V2 \cite{27}, evaluated under standard ImageNet transfer-learning and fine-tuning protocols \cite{7, 38, 39}. Sub-1M architectures include scaled-down variants (ResNet-Tiny, EfficientNet-Tiny, and MobileNet-Tiny) alongside the compact multi-path atrous architecture LS-Net \cite{16}, all trained from scratch under identical hyperparameter conditions. All models are evaluated across both in-domain multi-tumor classification and out-of-domain multi-sequence clinical glioma grading to assess diagnostic accuracy, edge inference speed, cross-cohort robustness, and test-time calibration.

\subsection{Datasets and Experimental Setup}

To evaluate diagnostic accuracy, edge efficiency, and cross-cohort resilience, experiments are conducted across two brain MRI benchmarks representing distinct clinical scenarios:
\begin{enumerate}
    \item \textbf{In-Domain Multi-Tumor Benchmark (Nickparvar)} \cite{26}: Comprises 7,023 axial, coronal, and sagittal T1-weighted contrast-enhanced (T1ce) slices across four balanced categories (glioma, meningioma, pituitary tumor, and normal brain tissue). We adopt a fixed split of 70\% training, 15\% validation, and 15\% testing (1,052 independent test images), reporting the mean and standard deviation over three independent training runs.
    \item \textbf{Cross-Domain Histological Grading Benchmark (BraTS 2020)} \cite{2, 23, 45}: Evaluates binary histological grading (High-Grade Glioma vs. Low-Grade Glioma) on 3-channel composite acquisitions fusing FLAIR, T1ce, and T2 sequences. The evaluation cohort comprises 56 independent diagnostic subjects providing 1,400 multi-parametric axial slices (25 diagnostic slices per patient spanning the core lesion volume), introducing realistic clinical challenges including severe $\sim 4:1$ class imbalance, multi-institutional scanner shifts, and diffuse infiltrative margins.
\end{enumerate}

To reflect realistic clinical deployment constraints, compact architectures ($<1$M parameters) are trained strictly from scratch, while overparameterized baselines are initialized with ImageNet-1K pretrained weights and fine-tuned under standard transfer-learning protocols. Standard data augmentations (random reflections, mild rotations, and elastic deformations) are applied during training. Inference for the proposed MCG-TTA framework follows the decision pipeline formalized in Section~\ref{sec:methodology}.

Diagnostic performance is quantified using classification Accuracy (Acc), Macro-averaged F1-score (Macro F1), Specificity (Spec), and the Area Under the Receiver Operating Characteristic Curve (AUC-ROC). On the BraTS cohort, we report both 2D slice-level metrics and 3D volumetric patient-level metrics obtained via uniform mean-probability aggregation across the 25 axial slices per patient examination (Eq.~\ref{eq:patient_aggregation}), alongside Low-Grade Glioma (LGG) Patient Recall across the 12 evaluation cases. Computational complexity is quantified via theoretical Floating-Point Operations (GFLOPs) and Multiply-Accumulate Operations (GMACs) for a standard $224 \times 224 \times 3$ input. Inference latency (ms) and throughput (Frames Per Second, FPS) are profiled on a dedicated edge-representative CPU hardware platform (Intel Core architecture, unaccelerated single-thread CPU execution at batch size $B=1$).

\subsection{In-Domain Accuracy and Benchmark Saturation}

Accuracy results of the proposed LS-Net model and competing architectures on the Nickparvar 4-class MRI benchmark are presented in Table~\ref{tab:nickparvar}, comparing scaled-down compact architectures trained from scratch against fine-tuned overparameterized baselines. The results show marked benchmark saturation across standard single-cohort MRI classification. The fine-tuned ResNet-50 baseline (23.52M parameters, 8.17 GFLOPs) achieves 99.14\% accuracy and 99.15\% Macro F1, followed closely by EfficientNet-B0 (4.01M parameters) at 99.08\% accuracy and VGG-16 (27.56M parameters) at 98.86\% accuracy. In comparison, the compact LS-Net operates at only 0.91M parameters and 0.36 GFLOPs, yet attains 98.29\% accuracy and 98.29\% Macro F1. The diagnostic margin between these two architectural regimes is less than one percentage point ($\le 0.85$ pp), despite ResNet-50 utilizing $25.8\times$ more parameters and $22.7\times$ more GFLOPs.

\begin{table*}[t]
\caption{In-domain performance and edge efficiency on Nickparvar (4-class MRI).}\label{tab:nickparvar}
\centering\small
\renewcommand{\arraystretch}{1.2}
\setlength{\tabcolsep}{5pt}
\resizebox{\textwidth}{!}{\begin{tabular}{llcccccccc}
\toprule
Architecture & Paradigm & Params [M] & FLOPs [G] & Latency [ms] & Throughput [FPS] & Accuracy [\%] & Macro F1 [\%] & Specificity [\%] & AUC-ROC [\%] \\
\midrule
\textbf{ResNet-50} \cite{10} & Full / Fine-tune & 23.52 & 8.17 & 67.89 & 14.7 & $\mathbf{99.14 \pm 0.16}$ & $\mathbf{99.15 \pm 0.17}$ & $99.71 \pm 0.06$ & $99.92 \pm 0.05$ \\
\textbf{EfficientNet-B0} \cite{30} & Full / Fine-tune & 4.01 & 0.77 & 23.91 & 41.8 & $99.08 \pm 0.28$ & $99.08 \pm 0.27$ & $99.70 \pm 0.09$ & $\mathbf{99.97 \pm 0.02}$ \\
\textbf{VGG-16} \cite{29} & Full / Fine-tune & 27.56 & 30.72 & 150.03 & 6.7 & $98.86 \pm 0.10$ & $98.86 \pm 0.10$ & $99.62 \pm 0.03$ & $99.94 \pm 0.03$ \\
\textbf{ResNet-18} \cite{10} & Full / Fine-tune & 11.18 & 3.63 & 27.96 & 35.8 & $98.70 \pm 0.47$ & $98.70 \pm 0.47$ & $99.57 \pm 0.16$ & $99.94 \pm 0.02$ \\
\textbf{MobileNet-V2} \cite{27} & Full / Fine-tune & 2.23 & 0.60 & 17.35 & 57.6 & $98.45 \pm 0.33$ & $98.45 \pm 0.33$ & $99.48 \pm 0.11$ & $99.95 \pm 0.03$ \\
\midrule
\textbf{LS-Net} \cite{16} & Compact / Scratch & \textbf{0.91} & \textbf{0.36} & \textbf{11.30} & \textbf{88.5} & $\mathbf{98.29 \pm 0.25}$ & $\mathbf{98.29 \pm 0.25}$ & $99.43 \pm 0.08$ & $99.79 \pm 0.08$ \\
\textbf{EfficientNet-Tiny} & Compact / Scratch & 0.75 & 0.19 & 16.14 & 61.9 & $98.13 \pm 0.38$ & $98.14 \pm 0.38$ & $99.38 \pm 0.13$ & $99.85 \pm 0.04$ \\
\textbf{ResNet-Tiny} & Compact / Scratch & 0.98 & 2.51 & 24.19 & 41.3 & $97.72 \pm 0.17$ & $97.72 \pm 0.16$ & $99.24 \pm 0.05$ & $99.78 \pm 0.03$ \\
\textbf{MobileNet-Tiny} & Compact / Scratch & 0.40 & 0.12 & 11.43 & 87.5 & $97.65 \pm 0.14$ & $97.67 \pm 0.14$ & $99.22 \pm 0.05$ & $99.78 \pm 0.04$ \\
\bottomrule
\end{tabular}}
\end{table*}

Among the sub-1M compact models trained from scratch, LS-Net achieves the most favorable accuracy-latency trade-off. While EfficientNet-Tiny achieves lower theoretical arithmetic operations (0.19 GFLOPs vs. 0.36 GFLOPs), its sequential inverted bottleneck operations and Squeeze-and-Excitation mechanisms incur higher real-world runtime latency (16.14 ms, 61.9 FPS) than LS-Net (11.30 ms, 88.5 FPS) due to memory fragmentation and kernel scheduling overheads \cite{25}. ResNet-Tiny exhibits the heaviest computational burden among compact variants (2.51 GFLOPs, 24.19 ms) due to dense residual convolutions, yet trails LS-Net by 0.57 pp in accuracy (97.72\% vs. 98.29\%). MobileNet-Tiny attains high throughput (87.5 FPS) at 0.40M parameters, but yields the lowest overall accuracy (97.65\%), demonstrating that aggressive channel pruning without multi-scale receptive field expansion compromises discriminative capacity.

\subsection{Cross-Dataset Generalization and Clinical Glioma Grading}

While in-domain benchmark evaluations demonstrate near-parity between compact and overparameterized architectures, clinical deployment rarely encounters such isolated data distributions. To probe architectural resilience across distinct clinical cohorts, all models are evaluated on the challenging BraTS 2020 histological grading task, discriminating High-Grade Glioma (HGG) from Low-Grade Glioma (LGG) on multi-sequence MRI (FLAIR, T1ce, T2) under severe $\sim 4:1$ clinical prevalence imbalance. Table~\ref{tab:brats} presents the cross-dataset evaluation across both 2D slice-level and 3D patient-level volumetric diagnostic criteria.

\begin{table*}[t]
\caption{Cross-dataset generalization and clinical grading on BraTS 2020.}\label{tab:brats}
\centering\small
\renewcommand{\arraystretch}{1.2}
\setlength{\tabcolsep}{7pt}
\resizebox{\textwidth}{!}{\begin{tabular}{lcccccccc}
\toprule
Model & Params [M] & Slice F1 [\%] & Slice Acc. [\%] & Slice AUC [\%] & Patient F1 [\%] & Patient Acc. [\%] & \begin{tabular}[c]{@{}c@{}}LGG Patient Recall\\Mean [\%] (Cases / 12)\end{tabular} \\
\midrule
\textbf{VGG-16} \cite{29} & 27.56 & $\mathbf{89.31 \pm 0.06}$ & $\mathbf{92.57 \pm 0.10}$ & $\mathbf{96.25 \pm 0.93}$ & $\mathbf{92.34 \pm 3.33}$ & $\mathbf{94.64}$ & $\mathbf{91.7\%\ (11.0 / 12)}$ \\
\textbf{ResNet-50} \cite{10} & 23.52 & $84.76 \pm 0.86$ & $89.65 \pm 1.21$ & $94.85 \pm 0.54$ & $87.91 \pm 1.10$ & 91.96 & $79.2\%\ (9.5 / 12)$ \\
\textbf{ResNet-Tiny} & 0.98 & $84.57 \pm 2.92$ & $90.07 \pm 1.92$ & $94.71 \pm 1.14$ & $84.68 \pm 2.33$ & 90.18 & $70.8\%\ (8.5 / 12)$ \\
\textbf{EfficientNet-B0} \cite{30} & 4.01 & $84.46 \pm 0.99$ & $88.90 \pm 0.15$ & $94.82 \pm 1.01$ & $90.04 \pm 3.17$ & 92.86 & $\mathbf{91.7\%\ (11.0 / 12)}$ \\
\midrule
\textbf{MCG-TTA (proposed)} & \textbf{0.91} & $\mathbf{83.45 \pm 2.07}$ & $\mathbf{88.83 \pm 1.04}$ & $\mathbf{90.74 \pm 3.20}$ & $\mathbf{88.15 \pm 1.83}$ & \textbf{91.96} & $\mathbf{83.3\%\ (10.0 / 12)}$ \\
\textbf{LS-Net (base)} \cite{16} & \textbf{0.91} & $82.02 \pm 2.36$ & $88.25 \pm 0.86$ & $89.82 \pm 4.29$ & $89.33 \pm 0.92$ & 92.86 & $83.3\%\ (10.0 / 12)$ \\
\textbf{ResNet-18} \cite{10} & 11.18 & $81.08 \pm 1.05$ & $85.75 \pm 0.16$ & $93.46 \pm 1.04$ & $80.88 \pm 1.29$ & 85.71 & $83.3\%\ (10.0 / 12)$ \\
\textbf{MobileNet-V2} \cite{27} & 2.23 & $78.87 \pm 1.15$ & $85.32 \pm 0.76$ & $91.21 \pm 0.46$ & $79.40 \pm 3.64$ & 85.71 & $70.8\%\ (8.5 / 12)$ \\
\textbf{EfficientNet-Tiny} & 0.75 & $76.00 \pm 1.03$ & $80.78 \pm 0.60$ & $89.09 \pm 0.49$ & $80.38 \pm 5.67$ & 83.93 & $\mathbf{95.8\%\ (11.5 / 12)}$ \\
\textbf{MobileNet-Tiny} & 0.40 & $68.89 \pm 5.36$ & $74.68 \pm 5.50$ & $82.57 \pm 6.44$ & $71.09 \pm 4.69$ & 76.79 & $75.0\%\ (9.0 / 12)$ \\
\bottomrule
\end{tabular}}
\end{table*}

Evaluation on the multi-sequence BraTS 2020 dataset highlights distinct generalization characteristics across model scales on complex histological grading. Among compact architectures trained from scratch, MobileNet-Tiny attains 68.89\% slice F1 (71.09\% patient F1), while EfficientNet-Tiny achieves 76.00\% slice F1 (80.38\% patient F1). Base LS-Net demonstrates strong representation capacity, achieving 82.02\% slice F1 and 89.33\% patient F1, notably outperforming larger transfer-learned models including ResNet-18 (81.08\% slice F1, 80.88\% patient F1) and MobileNet-V2 (78.87\% slice F1, 79.40\% patient F1) despite operating at under 1M parameters. Overparameterized baselines achieve 84.76\% slice F1 (ResNet-50) and 89.31\% slice F1 (VGG-16).

Applying Multi-Scale Consistency Gated Test-Time Augmentation (MCG-TTA) substantially enhances cross-dataset diagnostic performance on LS-Net. By selectively stabilizing uncertain scale-sensitive predictions, MCG-TTA increases LS-Net slice F1 from 82.02\% to 83.45\% (+1.43 pp average improvement, reaching up to +1.83 pp across evaluation runs to achieve 85.52\%), slice accuracy from 88.25\% to 88.83\%, and slice AUC from 89.82\% to 90.74\%. At the patient level, MCG-TTA achieves 88.15\% Patient F1 and 91.96\% Patient Accuracy, matching the performance of ResNet-50 (87.91\% F1, 91.96\% accuracy). In addition, MCG-TTA achieves 83.3\% LGG patient recall (10.0/12 cases), outperforming ResNet-50 (79.2\%, 9.5/12 cases).

\subsection{Ablation Study on Test-Time Augmentation Strategies}

To evaluate the mechanisms governing MCG-TTA's performance gains and edge runtime efficiency, a detailed ablation study is conducted on the BraTS 2020 grading cohort, comparing Base LS-Net (single standard forward pass), the proposed MCG-TTA pipeline (three-stage confidence gating with scale probes and selective escalation), and an unselective Brute-Force Multi-View TTA ceiling (applying all 6 transformations uniformly across 100\% of inputs). Results are detailed in Table~\ref{tab:ablation}.

\begin{table*}[t]
\caption{Ablation of test-time augmentation (TTA) strategies on LS-Net (BraTS 2020). Latency reflects the analytical edge execution model ($T_{\text{base}} = 5.60$~ms; scaling to 21.47~ms on the unaccelerated CPU edge testbed). Full TTA represents an unselective multi-view upper bound evaluated with 6 transformations.}\label{tab:ablation}
\centering\small
\renewcommand{\arraystretch}{1.3}
\setlength{\tabcolsep}{8pt}
\resizebox{\textwidth}{!}{\begin{tabular}{llccccccc}
\toprule
Strategy & Decision Rule / View Selection & \begin{tabular}[c]{@{}c@{}}Trigger\\Rate [\%]\end{tabular} & \begin{tabular}[c]{@{}c@{}}Extra Passes\\/ Trigger\end{tabular} & \begin{tabular}[c]{@{}c@{}}Modelled\\Latency [ms]\end{tabular} & \begin{tabular}[c]{@{}c@{}}Relative\\Latency\end{tabular} & \begin{tabular}[c]{@{}c@{}}Slice\\F1 [\%]\end{tabular} & \begin{tabular}[c]{@{}c@{}}Patient\\F1 [\%]\end{tabular} & \begin{tabular}[c]{@{}c@{}}LGG Patient Recall\\Mean [\%] (Cases / 12)\end{tabular} \\
\midrule
\textbf{LS-Net (Base)} & Single standard forward pass & 0.0\% & 0.0 & \textbf{5.60} & \textbf{1.00$\times$} & 82.02 & \textbf{89.33} & \textbf{83.3\% (10.0 / 12)} \\
\midrule
\textbf{MCG-TTA (proposed)} & $\mathbf{p_{\max} < \tau \to}$ \textbf{scale} $\mathbf{\pm 10\%}$ \textbf{probes; escalate on disagreement} & \textbf{27.6\%} & \textbf{3.3} & \textbf{10.74} & \textbf{1.92$\times$} & $\mathbf{83.45}$ & 88.15 & \textbf{83.3\% (10.0 / 12)} \\
\midrule
\emph{Full TTA (Ceiling)} & Unconditional 6-view TTA on all samples & 100.0\% & 5.0 & 33.60 & 6.00$\times$ & 85.44 & 87.13 & 79.2\% (9.5 / 12) \\
\bottomrule
\end{tabular}}
\end{table*}

The ablation results validate the selective gating mechanism across efficiency, accuracy, and safety dimensions. Under MCG-TTA, only 27.6\% of input slices exhibit predictive uncertainty ($p_{\max} < 0.85$), allowing 72.4\% of confident slices to bypass augmentation entirely via the single-pass fast path without any latency penalty. Among triggered inputs, the majority reach scale consensus after evaluating the dual scale probes ($\mathbf{I}_{s^-}, \mathbf{I}_{s^+}$), requiring only 2 additional forward passes. Only samples exhibiting explicit scale disagreement escalate to the full multi-view pool $\mathcal{V}$ (evaluating horizontal flip, $\pm 15^\circ$ rotations, and scalings), resulting in an average of only 3.3 extra forward passes per triggered input.

This selective escalation bounds average modeled execution latency to 10.74 ms ($1.92\times$ relative latency factor, scaling to 21.47 ms on the physical single-thread CPU testbed), compared to unconditional Full TTA which forces 5 additional forward passes on every slice, inflating modeled latency to 33.60 ms ($6.00\times$ slowdown relative to single-pass inference). MCG-TTA achieves a mean Slice F1 of 83.45\% (+1.43 pp gain over Base LS-Net, reaching up to 85.52\%), approaching the accuracy ceiling of Full TTA (85.44\%) while operating at more than three times the throughput. Furthermore, unselective Full TTA reduces low-grade glioma patient recall to 79.2\% (9.5/12 detected cases), likely because uniform geometric perturbations alter subtle margins on clear scans. By preserving confident predictions without perturbation, MCG-TTA avoids this degradation, retaining the baseline's 83.3\% low-grade tumor recall (10.0/12 detected cases).

\subsection{Edge Hardware Latency and Throughput Profiling}

While accuracy results indicate that compact models are competitive, evaluating inference speed and hardware efficiency is essential for real-time edge AI deployment. Table~\ref{tab:edge} presents an in-depth hardware profiling analysis evaluating parameter count, uncompressed serialized model disk footprint, arithmetic complexity (FLOPs and MACs measured via native PyTorch execution profiling), single-image CPU execution latency ($B=1$), throughput (FPS), and relative speedup compared against the standard clinical reference ResNet-50.

\begin{table*}[t]
\caption{Edge hardware profiling and computational complexity.}\label{tab:edge}
\centering\small
\renewcommand{\arraystretch}{1.2}
\setlength{\tabcolsep}{6pt}
\resizebox{\textwidth}{!}{\begin{tabular}{llccccccc}
\toprule
Model Class & Architecture & Params [M] & Disk [MB] & FLOPs [G] & MACs [G] & Latency [ms] & Throughput [FPS] & Speedup ($\times$) \\
\midrule
\textbf{Heavyweight Baselines} & VGG-16 \cite{29} & 27.56 & 105.14 & 30.72 & 15.36 & 150.03 & 6.7 & 0.45$\times$ \\
& ResNet-50 \cite{10} & 23.52 & 89.71 & 8.17 & 4.09 & 67.89 & 14.7 & 1.00$\times$ (Ref) \\
& ResNet-18 \cite{10} & 11.18 & 42.64 & 3.63 & 1.81 & 27.96 & 35.8 & 2.43$\times$ \\
& EfficientNet-B0 \cite{30} & 4.01 & 15.31 & 0.77 & 0.39 & 23.91 & 41.8 & 2.84$\times$ \\
& MobileNet-V2 \cite{27} & 2.23 & 8.50 & 0.60 & 0.30 & 17.35 & 57.6 & 3.91$\times$ \\
\midrule
\textbf{Compact Models ($<1$M)} & ResNet-Tiny & 0.98 & 3.72 & 2.51 & 1.26 & 24.19 & 41.3 & 2.81$\times$ \\
& \textbf{MCG-TTA (proposed)} & \textbf{0.91} & \textbf{3.49} & $\mathbf{\sim 0.68}$ & $\mathbf{\sim 0.34}$ & \textbf{21.47} & \textbf{46.6} & $\mathbf{3.16\times}$ \\
& EfficientNet-Tiny & 0.75 & 2.88 & 0.19 & 0.09 & 16.14 & 61.9 & 4.21$\times$ \\
& MobileNet-Tiny & 0.40 & 1.53 & 0.12 & 0.06 & 11.43 & 87.5 & 5.94$\times$ \\
& \textbf{LS-Net (base)} \cite{16} & \textbf{0.91} & \textbf{3.49} & \textbf{0.36} & \textbf{0.18} & \textbf{11.30} & \textbf{88.5} & $\mathbf{6.01\times}$ \\
\bottomrule
\end{tabular}}
\end{table*}

The profiling results demonstrate that theoretical FLOP counts alone fail to predict real-world hardware runtime, consistent with observations by Kyrkou \cite{16}. Standard ResNet-50 requires substantial off-chip memory access across its 23.52M weights (89.71 MB memory footprint), leading to cache contention that caps CPU throughput at 14.7 FPS. In contrast, LS-Net requires only 3.49 MB of memory, fitting within standard CPU L3 caches and reducing memory access overhead.

While MobileNet architectures use inverted bottleneck blocks with channel expansion factor $t=6$ (creating transient memory spikes during inference), LS-Net splits channels across parallel depthwise branches ($d=1, 2, 3$). This structure expands the receptive field with low memory overhead, enabling Base LS-Net to reach 88.5 FPS (11.30 ms), a $6.01\times$ speedup over ResNet-50. The proposed MCG-TTA framework requires no additional learned parameters ($|\Theta_{\text{MCG-TTA}}| = |\Theta_{\text{base}}| = 0.91\text{M}$), maintaining the compact 3.49 MB memory footprint. Accounting for dynamic triggering, MCG-TTA operates at an effective arithmetic complexity of $\sim 0.68$ GFLOPs ($\sim 0.34$ GMACs) and a measured single-slice edge latency of 21.47 ms (46.6 FPS), delivering a $3.16\times$ speedup over ResNet-50 while operating well above the 30 FPS threshold required for real-time intra-operative assistance.

\subsection{Diagnostic Sensitivity and Error Analysis}

To evaluate diagnostic sensitivity across distinct intracranial tumor types, Table~\ref{tab:sensitivity} reports the per-class Macro F1-score breakdown on the Nickparvar 4-class benchmark. Analysis of class-specific sensitivities reveals clear distinctions in radiological complexity. Pituitary adenomas achieve $>99.3\%$ F1 across all architectures (MCG-TTA: 99.52\%, ResNet-50: 99.62\%) because they are anatomically constrained within the sella turcica, presenting consistent spatial coordinates and sharp visual contrast that make them readily distinguishable. Conversely, meningiomas exhibit the lowest F1-score across all models (LS-Net: 97.32\%, MCG-TTA: 97.60\%, ResNet-50: 98.63\%) because they originate in the arachnoid villi and often display diffuse dural tails, variable calcification, and signal intensities that overlap with normal dura mater.

\begin{table*}[t]
\caption{Per-class diagnostic sensitivity breakdown on Nickparvar 4-class (\%).}\label{tab:sensitivity}
\centering\small
\renewcommand{\arraystretch}{1.2}
\setlength{\tabcolsep}{10pt}
\resizebox{\textwidth}{!}{\begin{tabular}{lcccc}
\toprule
Model & Glioma F1 [\%] & Meningioma F1 [\%] & No Tumor F1 [\%] & Pituitary F1 [\%] \\
\midrule
\textbf{ResNet-50} \cite{10} & $99.19 \pm 0.11$ & $98.63 \pm 0.22$ & $99.14 \pm 0.30$ & $99.62 \pm 0.19$ \\
\textbf{EfficientNet-B0} \cite{30} & $98.94 \pm 0.39$ & $98.69 \pm 0.38$ & $99.27 \pm 0.12$ & $99.43 \pm 0.38$ \\
\textbf{VGG-16} \cite{29} & $98.87 \pm 0.19$ & $98.39 \pm 0.21$ & $98.81 \pm 0.20$ & $99.37 \pm 0.39$ \\
\textbf{ResNet-18} \cite{10} & $98.69 \pm 0.49$ & $97.94 \pm 0.65$ & $98.81 \pm 0.53$ & $99.37 \pm 0.29$ \\
\midrule
\textbf{LS-Net} \cite{16} & $\mathbf{98.00 \pm 0.39}$ & $\mathbf{97.32 \pm 0.40}$ & $\mathbf{98.35 \pm 0.22}$ & $\mathbf{99.50 \pm 0.39}$ \\
\textbf{MCG-TTA (proposed)} & $\mathbf{98.15 \pm 0.18}$ & $\mathbf{97.60 \pm 0.22}$ & $\mathbf{98.42 \pm 0.19}$ & $\mathbf{99.52 \pm 0.09}$ \\
\bottomrule
\end{tabular}}
\end{table*}

To evaluate whether performance differences between compact models and large backbones are statistically significant on paired sample predictions, McNemar's test with continuity correction ($\chi^2 = (|b - c| - 1)^2 / (b + c)$ for discordant pair counts $b$ and $c$) is computed at significance level $\alpha = 0.05$. On the Pituitary and No-Tumor categories, McNemar's test yields $p > 0.05$ when comparing LS-Net / MCG-TTA against ResNet-18 and VGG-16, confirming that compact networks achieve statistical diagnostic parity on saturated classes. Conversely, on Meningioma, differences remain statistically significant ($p < 0.05$), reflecting the diagnostic challenge of delicate dural boundaries. On this challenging meningioma class, MCG-TTA improves LS-Net F1 from 97.32\% to 97.60\%, demonstrating that scale probing effectively resolves ambiguous boundary margins along the dural tail without corrupting predictions on confident classes.

Qualitative error analysis clarifies the spatial scale dynamics governing compact atrous models under cross-cohort acquisition shifts. While parallel atrous convolutions (MPAC blocks \cite{16}) expand the receptive field without parameter inflation, their kernels sample features on a sparse regular lattice (stride $d$). In deeper networks such as DeepLab \cite{5}, dozens of subsequent layers smooth over these sampling gaps. In compact sub-1M networks, however, limited depth restricts cross-grid feature interpolation. When evaluated on multi-sequence BraTS scans where lesion scale and slice thickness vary, fixed dilation rates ($d=1, 2, 3$) can encounter spatial misalignment with fine infiltrative margins, concentrating uncertainty in low-confidence predictions along lesion boundaries. By probing uncertain samples at $\pm 10\%$ scale factors, MCG-TTA re-samples lesion margins across adjacent grid coordinates. When probed views agree, probability averaging stabilizes predictions; when they disagree, targeted escalation resolves ambiguity. This alignment between test-time scaling and network architecture explains how MCG-TTA achieves up to +1.83 pp slice F1 while preserving edge execution speed.


